\documentclass[11pt,a4paper]{article}

\usepackage{fontspec}
\setmainfont{DejaVu Serif}[
  BoldFont={DejaVu Serif Bold},
  ItalicFont={DejaVu Serif},
  ItalicFeatures={FakeSlant=0.18},
  BoldItalicFont={DejaVu Serif Bold},
  BoldItalicFeatures={FakeSlant=0.18}
]
\setsansfont{DejaVu Sans}
\usepackage{polyglossia}
\setmainlanguage{english}
\usepackage[a4paper,margin=20.5mm]{geometry}
\usepackage{microtype}
\usepackage{setspace}
\setstretch{1.012}
\usepackage{amsmath,amssymb,amsthm,mathtools,bm}
\usepackage{booktabs,array,tabularx}
\usepackage{enumitem}
\usepackage{xcolor}
\usepackage{tikz}
\usetikzlibrary{arrows.meta,positioning,calc}
\usepackage[unicode,hidelinks]{hyperref}
\usepackage[nameinlink,noabbrev]{cleveref}
\emergencystretch=3em
\allowdisplaybreaks

\hypersetup{
  pdftitle={Quantum Chemistry as a Typed Inference Stack under Observation Maps},
  pdfauthor={Stas, Independent Research Program},
  pdfsubject={Many-electron states, reduced density matrices, orbital gauge, nuclear geometry, normal modes, spectra, and observation channels},
  pdfkeywords={quantum chemistry, reduced density matrix, orbital gauge, Born-Oppenheimer approximation, normal modes, spectroscopy, observation geometry}
}

\definecolor{obsblue}{RGB}{37,83,138}
\definecolor{obsred}{RGB}{150,45,45}
\definecolor{obsgreen}{RGB}{30,115,78}
\definecolor{obsviolet}{RGB}{94,64,145}
\definecolor{lightblue}{RGB}{239,246,253}
\definecolor{lightred}{RGB}{253,242,242}
\definecolor{lightgreen}{RGB}{239,249,244}
\definecolor{lightviolet}{RGB}{246,242,252}

\theoremstyle{definition}
\newtheorem{definition}{Definition}[section]
\newtheorem{model}[definition]{Model}
\theoremstyle{plain}
\newtheorem{theorem}[definition]{Theorem}
\newtheorem{proposition}[definition]{Proposition}
\newtheorem{lemma}[definition]{Lemma}
\newtheorem{corollary}[definition]{Corollary}
\theoremstyle{remark}
\newtheorem{remark}[definition]{Remark}
\newtheorem{warning}[definition]{Boundary}

\newcommand{\R}{\mathbb{R}}
\newcommand{\C}{\mathbb{C}}
\newcommand{\dd}{\mathop{}\!\mathrm d}
\newcommand{\norm}[1]{\left\lVert #1\right\rVert}
\newcommand{\abs}[1]{\left\lvert #1\right\rvert}
\newcommand{\pair}[2]{\left\langle #1,#2\right\rangle}
\newcommand{\Hcal}{\mathcal H}
\newcommand{\Qcal}{\mathcal Q}
\newcommand{\Scal}{\mathcal S}
\newcommand{\Acal}{\mathcal A}
\newcommand{\Bcal}{\mathcal B}
\newcommand{\I}{\mathrm i}
\newcommand{\Tr}{\operatorname{Tr}}
\newcommand{\rank}{\operatorname{rank}}
\newcommand{\diag}{\operatorname{diag}}
\newcommand{\Span}{\operatorname{span}}
\newcommand{\SE}{\operatorname{SE}}
\newcommand{\ket}[1]{\left\lvert #1\right\rangle}
\newcommand{\bra}[1]{\left\langle #1\right\rvert}
\newcolumntype{L}[1]{>{\raggedright\arraybackslash}p{#1}}
\newcolumntype{Y}{>{\raggedright\arraybackslash}X}

\title{\bfseries Quantum Chemistry as a Typed Inference Stack\\
\large Many-electron states, reduced data, orbital gauge, nuclear geometry, and spectra}
\author{Stas\\
\small Independent Research Program; ORCID: 0009-0000-2294-705X}
\date{Strict GO Core reference revision v1.1 --- 28 July 2026}

\begin{document}
\maketitle

\begin{center}
\fcolorbox{obsblue}{lightblue}{%
\parbox{0.92\linewidth}{\small
\textbf{Status.} This document supersedes the bilingual working note
v1.0 and is an application-layer reference migrated to GO Core v0.2.
It does not introduce a new electronic-structure method, density
functional, orbital interpretation, bond definition, or spectroscopic
inversion theorem. Its internal contribution is a typed separation of
the exact electron-nuclear state, clamped-nuclei electronic models,
reduced density data, representation gauge, nuclear reduction,
ideal observables, and finite-resolution measurement.}}
\end{center}

\begin{abstract}
Quantum chemistry is not described by a single composable chain
\(\Psi\to n\to E(R)\to R_\ast\to\) chemical behavior.  Those arrows
have different domains, hypotheses, and information content.  This
reference replaces that chain by a typed graph.  The exact
nonrelativistic molecular state depends on electronic and nuclear
coordinates.  A clamped-nuclei Hamiltonian family
\(H_{\rm e}(R)\) generates potential-energy surfaces only after a
Born--Oppenheimer reduction.  For an \(N\)-electron density operator
\(\Gamma_N\), the one-body reduced density matrix
\(\gamma=N\Tr_{2\ldots N}\Gamma_N\) satisfies
\(\Tr\gamma=N\), and the spatial density integrates to \(N\), but this
reduction is non-injective.  Occupied-orbital rotations leave a Slater
determinant ray and its one-body projector invariant, so localized
orbitals, hybridizations, and active-space populations require an
explicit representation protocol.  At a stationary nuclear geometry,
the mass-weighted Hessian has eigenvalues
\(\lambda_k=\omega_k^2\), not frequencies; five rigid zero modes for a
linear molecule or six for a nonlinear molecule must be removed in
the isolated, unconstrained case.  An ideal transition spectrum
depends on energies, populations, a transition operator, and a line
profile, while measured data additionally depend on the instrument
channel and noise.  The resulting formalism preserves the useful
observation-geometry interpretation but forbids claims of canonical
orbital observability, unique bond projection, exact degeneracy from
finite resolution, or chemistry determined by energy differences
alone.
\end{abstract}

\tableofcontents

\section{Scope, state spaces, and unit firewall}

Let \(N_{\rm e}\) electrons and \(N_{\rm n}\) nuclei be described in a
declared inertial frame.  Electronic spin coordinates are denoted by
\(x_i=(\bm r_i,\sigma_i)\), and nuclear Cartesian coordinates by
\(R=(\bm R_1,\ldots,\bm R_{N_{\rm n}})\).  Ignoring relativity,
external fields, nuclear spin, and finite nuclear size, the
nonrelativistic Coulomb Hamiltonian in atomic units is
\begin{equation}\label{eq:full-hamiltonian}
\begin{split}
\widehat H_{\rm mol}
={}&-\frac12\sum_{i=1}^{N_{\rm e}}\nabla_i^2
-\sum_{A=1}^{N_{\rm n}}\frac{1}{2M_A}\nabla_A^2
-\sum_{i,A}\frac{Z_A}{r_{iA}}\\
&+\sum_{i<j}\frac{1}{r_{ij}}
+\sum_{A<B}\frac{Z_AZ_B}{R_{AB}} .
\end{split}
\end{equation}
The nuclear masses \(M_A\) are expressed in electron-mass units.  A
full stationary state is
\begin{equation}\label{eq:full-state}
\Psi=\Psi(x_1,\ldots,x_{N_{\rm e}};R_1,\ldots,R_{N_{\rm n}}),
\qquad
\widehat H_{\rm mol}\Psi=E\Psi.
\end{equation}
Writing only electronic arguments after \cref{eq:full-hamiltonian}
silently changes the problem and is therefore prohibited.

\begin{warning}[Unit context]
Atomic units suppress \(m_{\rm e},e,\hbar,4\pi\epsilon_0\), but do not
erase physical dimensions.  A numerical atomic-unit expression may
enter an SI formula only through an explicit conversion using the
declared values of \(E_{\rm h},a_0,m_{\rm e},e,\hbar\).  No addition or
comparison across SI and atomic contexts is implicit in this
reference.
\end{warning}

For pure isolated electronic states, the state space is the
antisymmetric sector
\[
\Hcal_{\rm e}^{(N)}
=\bigwedge\nolimits^N
L^2(\R^3;\C^2).
\]
Mixed or open-system electronic states are positive trace-one
operators \(\Gamma_N\) on this sector.  A wavefunction is therefore
not the universal state object; it is a pure-state representative.

\section{Exact molecule and Born--Oppenheimer reduction}

At fixed nuclear geometry \(R\), define the clamped-nuclei electronic
Hamiltonian
\begin{equation}\label{eq:clamped-hamiltonian}
\widehat H_{\rm e}(R)
=\widehat T_{\rm e}+\widehat V_{\rm ee}
+\widehat V_{\rm en}(R)+V_{\rm nn}(R).
\end{equation}
This document includes the scalar nuclear repulsion
\(V_{\rm nn}(R)\) in \(\widehat H_{\rm e}(R)\); a convention that
excludes it must add it back before calling an eigenvalue a
potential-energy surface.

The parametric eigenproblem is
\begin{equation}\label{eq:electronic-eigenproblem}
\widehat H_{\rm e}(R)\psi_k(x;R)=E_k(R)\psi_k(x;R).
\end{equation}
The map \(R\mapsto E_k(R)\) is not obtained by taking a marginal of the
exact molecular wavefunction.  It is generated from the Hamiltonian
family after choosing a clamped-nuclei model and an electronic state
label.  A Born--Huang expansion has the form
\begin{equation}\label{eq:born-huang}
\Psi(x,R)=\sum_k\chi_k(R)\psi_k(x;R).
\end{equation}
Substitution into the full equation yields coupled nuclear equations
containing derivative couplings.  A one-surface Born--Oppenheimer
model is a reduction obtained by neglecting specified off-diagonal
terms; it is not a lossless frame transform and it is not uniformly
valid at a closing internal electronic gap.

\begin{warning}[Conical-intersection inheritance]
Near a conical intersection, the individual rank-one adiabatic
projectors need not extend through the seam even when the retained
two-state cluster remains regular.  Branching coordinates, gauge,
Berry holonomy, derivative coupling, dynamics, and finite-resolution
identifiability are governed by the separate P8 conical-intersection
contract.  The present document does not replace that contract with
the phrase ``observation caustic.''
\end{warning}

\section{Reduced density data and information loss}

\begin{definition}[One-body reduced density matrix]
For a normalized \(N\)-electron state \(\Gamma_N\), define
\begin{equation}\label{eq:one-rdm}
\gamma^{(1)}
=N\,\Tr_{2,\ldots,N}\Gamma_N .
\end{equation}
Its spatial density is
\begin{equation}\label{eq:density}
n(\bm r)
=\sum_{\sigma}
\langle \bm r\sigma|\gamma^{(1)}|\bm r\sigma\rangle .
\end{equation}
\end{definition}

\begin{proposition}[Normalization and one-body observables]
\label{prop:rdm-normalization}
The reduction in \cref{eq:one-rdm} obeys
\begin{equation}\label{eq:rdm-normalization}
\Tr\gamma^{(1)}=N,
\qquad
\int_{\R^3}n(\bm r)\,\dd^3r=N.
\end{equation}
For every one-body observable
\(\widehat A=\sum_{i=1}^N a(i)\),
\begin{equation}\label{eq:one-body-expectation}
\Tr(\Gamma_N\widehat A)=\Tr(\gamma^{(1)}a).
\end{equation}
\end{proposition}

\begin{proof}
The partial trace preserves the trace, so
\(\Tr\gamma^{(1)}=N\Tr\Gamma_N=N\).  Resolving the one-particle trace
in the position-spin basis gives the density integral.  Permutation
symmetry of \(\Gamma_N\) makes all \(N\) terms in \(\widehat A\)
equal after tracing the other particles, which yields
\cref{eq:one-body-expectation}.
\end{proof}

\begin{proposition}[The one-body reduction is non-injective]
\label{prop:rdm-noninjective}
Let \(\ket{ij}=a_i^\dagger a_j^\dagger\ket0\) in
\(\bigwedge^2\C^4\), and define
\[
\ket{\Psi_\pm}=\frac{\ket{12}\pm\ket{34}}{\sqrt2}.
\]
Both states have
\(\gamma^{(1)}=\tfrac12 I_4\), but the two-body coherence
\[
\widehat B=\ket{12}\!\bra{34}+\ket{34}\!\bra{12}
\]
has expectation \(+1\) in \(\Psi_+\) and \(-1\) in \(\Psi_-\).
\end{proposition}

\begin{proof}
Each spin orbital occurs with probability \(1/2\), and all one-body
cross terms vanish because the determinants differ by two orbitals.
The operator \(\widehat B\) exchanges the two determinants, producing
the stated signs.
\end{proof}

Thus neither \(\Gamma_N\mapsto\gamma^{(1)}\) nor
\(\Gamma_N\mapsto n\) is generally invertible.  Electron correlation
is not defined as an information-theoretic defect of this map.  It is
method- and Hamiltonian-specific unless a precise correlation
functional or entanglement quantity is declared.

\begin{center}
\begin{tikzpicture}[
  node distance=6mm and 13mm,
  box/.style={draw=obsblue,rounded corners,fill=lightblue,
    align=center,inner sep=4pt,font=\scriptsize},
  modelbox/.style={draw=obsviolet,rounded corners,fill=lightviolet,
    align=center,inner sep=4pt,font=\scriptsize},
  obsbox/.style={draw=obsgreen,rounded corners,fill=lightgreen,
    align=center,inner sep=4pt,font=\scriptsize},
  arr/.style={-{Latex[length=2mm]},thick,draw=obsblue},
  marr/.style={-{Latex[length=2mm]},thick,draw=obsviolet},
  oarr/.style={-{Latex[length=2mm]},thick,draw=obsgreen},
  lab/.style={font=\tiny,fill=white,inner sep=1pt}
]
\node[box] (state) {\(\Gamma_N\)\\many-body state};
\node[box,right=of state] (rdm) {\(\gamma^{(1)}\)\\one-body reduction};
\node[box,right=of rdm] (density) {\(n(\bm r)\)\\density};
\draw[arr] (state) -- node[above,lab] {partial trace} (rdm);
\draw[arr] (rdm) -- node[above,lab] {diagonal} (density);

\node[modelbox,below=of state] (family) {\(H_{\rm e}(R)\)\\model family};
\node[modelbox,right=of family] (pes) {\(E_k(R)\)\\PES branch};
\node[modelbox,right=of pes] (geom) {\([R_\ast]\)\\geometry class};
\draw[marr] (family) -- node[above,lab] {eigensolve} (pes);
\draw[marr] (pes) -- node[above,lab] {optimize} (geom);

\node[obsbox,below=of family] (ideal) {\((\Gamma,A)\)\\state and probe};
\node[obsbox,right=of ideal] (spectrum) {\(S_A(\omega)\)\\ideal signal};
\node[obsbox,right=of spectrum] (data) {\(Y\)\\instrument data};
\draw[oarr] (ideal) -- node[above,lab] {Born rule} (spectrum);
\draw[oarr] (spectrum) -- node[above,lab] {channel} (data);
\end{tikzpicture}
\end{center}

The three rows are related by a declared scientific workflow, but
they are not a single automatic composition.

\section{Finite bases, Hartree--Fock, and orbital gauge}

Let \(\{\chi_\mu\}_{\mu=1}^K\) be a linearly independent, generally
nonorthogonal one-particle basis with overlap matrix
\[
S_{\mu\nu}=\langle\chi_\mu,\chi_\nu\rangle,
\qquad S=S^\dagger>0.
\]
Roothaan--Hall equations have the generalized form
\begin{equation}\label{eq:roothaan}
FC=SC\varepsilon,
\qquad
C^\dagger SC=I.
\end{equation}
Writing \(FC=C\varepsilon\) without either an orthonormal basis or the
overlap matrix changes the discrete problem.  Small eigenvalues of
\(S\) also require a declared linear-dependence threshold.

For occupied columns \(C_{\rm o}\), define the AO density matrix
\begin{equation}\label{eq:ao-density}
P=C_{\rm o}C_{\rm o}^\dagger,
\qquad
\Tr(PS)=N_{\rm occ}
\end{equation}
for one electron per occupied spin orbital.  Closed-shell spatial
orbital conventions add their declared occupation factor.

\begin{proposition}[Occupied-orbital gauge]
\label{prop:occupied-gauge}
If \(U\in U(N_{\rm occ})\), then
\[
C_{\rm o}'=C_{\rm o}U
\quad\Longrightarrow\quad
P'=P.
\]
The associated normalized Slater determinant changes only by
\(\det U\), hence represents the same quantum ray.
\end{proposition}

\begin{proof}
\(P'=C_{\rm o}UU^\dagger C_{\rm o}^\dagger=P\).  Multilinearity of
the exterior product gives
\(\phi_1'\wedge\cdots\wedge\phi_{N_{\rm occ}}'
=(\det U)\phi_1\wedge\cdots\wedge\phi_{N_{\rm occ}}\), and
\(\abs{\det U}=1\).
\end{proof}

Consequently, localized occupied orbitals, hybrid orbitals, and
orbital-by-orbital bond pictures are representation choices unless
the localization or decomposition functional is part of the
protocol.  Canonical orbitals diagonalize a chosen effective
operator; that gauge choice does not turn every orbital or orbital
energy into a many-body observable.

Hartree--Fock is the variational problem
\begin{equation}\label{eq:hf-variational}
E_{\rm HF}
=\inf_{\Phi\ {\rm Slater},\,\norm{\Phi}=1}
\langle\Phi|\widehat H_{\rm e}|\Phi\rangle .
\end{equation}
The nonlinear SCF equations characterize stationary solutions.
An iteration may converge to a local stationary point, oscillate, or
break an exact symmetry; convergence alone does not certify the
global Hartree--Fock minimum.

For the same Hamiltonian and one-particle space,
\begin{equation}\label{eq:correlation-energy}
E_{\rm corr}=E_0-E_{\rm HF}\le0,
\end{equation}
where \(E_0\) is the exact or full-CI ground energy in the declared
space.  Cross-method correlation energies are not comparable unless
Hamiltonian, basis, frozen-core convention, geometry, and reference
state agree.

\section{Density-functional theory firewall}

The density map is a partial trace followed by a diagonal readout.
Ground-state density-functional theory adds a separate variational
statement.  For fixed particle number and electron-electron
interaction, a constrained-search functional can be written
\begin{equation}\label{eq:levy-search}
F[n]=\inf_{\Gamma\mapsto n}
\Tr\!\left[\Gamma(\widehat T+\widehat W)\right].
\end{equation}
The external-potential ground energy is then
\begin{equation}\label{eq:dft-variational}
E_v=\inf_n\left\{
F[n]+\int v_{\rm ext}(\bm r)n(\bm r)\,\dd^3r
\right\},
\end{equation}
over the declared admissible density class.

\begin{warning}[Hohenberg--Kohn scope]
In its original ground-state setting, the density determines the
external scalar potential up to an additive constant under fixed
particle number and interaction, subject to the theorem's
representability and degeneracy hypotheses.  This does not invert the
density map for an arbitrary excited, time-dependent, mixed, or
open-system state.  Degenerate and ensemble cases require their own
formulation.
\end{warning}

The Kohn--Sham construction introduces auxiliary orbitals with
\begin{equation}\label{eq:ks-density}
n(\bm r)=\sum_i f_i\abs{\varphi_i(\bm r)}^2,
\qquad
\left[-\frac12\nabla^2+v_{\rm eff}[n]\right]\varphi_i
=\varepsilon_i\varphi_i
\end{equation}
in atomic units.  The orbitals reproduce the target density within
the model.  General Kohn--Sham eigenvalue differences are not equal
by definition to interacting excitation energies, and approximate
exchange-correlation functionals introduce model error not encoded
by the density normalization.

\section{Nuclear geometry and the normal-mode theorem}

For an isolated molecule, absolute translation and rotation do not
define internal geometry.  Away from collisions and singular
configurations, an internal configuration is an equivalence class
\[
[R]\in
\Qcal_{\rm int}
=\bigl((\R^{3N_{\rm n}}\setminus\Delta)/\SE(3)\bigr)/
\Pi_{\rm identical},
\]
where the final quotient is included when identical-nucleus
permutations are physically identified.  A Cartesian minimizer is
therefore not unique before this quotient or a gauge-fixing
alignment.

Let \(R_\ast\) be a stationary geometry of a smooth PES,
\(H_R=\nabla_R^2E(R_\ast)\) its Cartesian Hessian, and
\(\mathsf M=\diag(M_1,M_1,M_1,\ldots)\).  The mass-weighted dynamical
matrix is
\begin{equation}\label{eq:mass-weighted-hessian}
\mathsf K
=\mathsf M^{-1/2}H_R\mathsf M^{-1/2}.
\end{equation}

\begin{theorem}[Harmonic normal modes]
\label{thm:normal-modes}
If \(\mathsf K u_k=\lambda_k u_k\), then the linearized nuclear
equation has modal solutions with
\begin{equation}\label{eq:frequency-square}
\lambda_k=\omega_k^2,
\qquad
\nu_k=\frac{\omega_k}{2\pi},
\qquad
\widetilde\nu_k=\frac{\omega_k}{2\pi c}.
\end{equation}
For an unconstrained isolated nonlinear molecule, six rigid
translation/rotation directions are removed; for a linear molecule,
five are removed.  A strict harmonic minimum has
\(\lambda_k>0\) on the remaining vibrational subspace.  A first-order
transition state has exactly one negative vibrational eigenvalue
there.
\end{theorem}

\begin{proof}
With displacement \(\delta R\), the linearized equation is
\(\mathsf M\delta\ddot R+H_R\delta R=0\).  Setting
\(q=\mathsf M^{1/2}\delta R\) gives
\(\ddot q+\mathsf Kq=0\).  Substitution
\(q(t)=u_ke^{\I\omega t}\) yields
\(\mathsf Ku_k=\omega^2u_k\).  Euclidean invariance makes
translations and infinitesimal rotations null directions at a
stationary geometry.  Rotation about the molecular axis displaces no
nucleus for a linear geometry, leaving five rather than six
independent rigid directions.
\end{proof}

\begin{warning}[Qualification]
Atoms, constraints, periodic cells, external fields, nonstationary
points, exact symmetries beyond rigid motion, and numerical
finite-difference Hessians change the null-space analysis.  A small
computed value is not automatically an exact zero mode; the threshold
must be scaled and reported.
\end{warning}

\section{Active spaces, valence, and bond interpretations}

A mathematically typed active-space reduction begins with a declared
one-particle projector \(P_{\Acal}\):
\begin{equation}\label{eq:active-rdm}
\gamma_{\Acal}
=P_{\Acal}\gamma^{(1)}P_{\Acal},
\qquad
N_{\Acal}=\Tr(\gamma^{(1)}P_{\Acal}).
\end{equation}
The projector may be generated by atomic valence functions, energy
windows, natural occupations, localized orbitals, or another stated
algorithm.  Different admissible choices need not give the same
\(N_{\Acal}\), oxidation assignment, or bond decomposition.

\begin{definition}[Interpretation map]
A population, partial charge, bond order, hybridization, or
``valence projection'' is an interpretation map
\[
\mathfrak I_\alpha:
(\Gamma,H,\Bcal,\Acal,\text{conventions})
\longmapsto z_\alpha
\]
with a named scheme \(\alpha\), basis \(\Bcal\), active space
\(\Acal\), and normalization.  It is not promoted to a
scheme-independent observable without an invariance theorem.
\end{definition}

This distinction does not make such quantities useless.  It makes
comparisons conditional: trends are meaningful within a controlled
protocol, while equality across unrelated partition schemes is not
assumed.

\section{Ideal spectra and finite-resolution data}

Let \(\{\ket i,E_i\}\) be stationary states, \(p_i\) initial
populations, and \(\widehat A\) the probe-dependent transition
operator.  With a normalized line profile
\(\int L_\Gamma(\omega)\dd\omega=1\), define
\begin{equation}\label{eq:ideal-spectrum}
S_A(\omega)
=\sum_{i,f}p_i
\abs{\langle f|\widehat A|i\rangle}^2
L_\Gamma\!\left(
\omega-\frac{E_f-E_i}{\hbar}
\right).
\end{equation}
The observed record is a separate channel, schematically
\begin{equation}\label{eq:instrument-channel}
Y
=\mathcal D_\epsilon
\mathcal C_{\rm inst}[S_A]+\eta,
\end{equation}
where \(\mathcal C_{\rm inst}\) includes response and convolution,
\(\mathcal D_\epsilon\) sampling or quantization, and \(\eta\) a
declared stochastic error process.

Energy differences determine line locations only after a dynamical
model identifies the relevant states.  Intensities also require
populations and matrix elements.  Broadening, unresolved blends,
calibration, baseline, and selection bias enter the inference from
\(Y\) back to a molecular model.

If a symmetry product contains the totally symmetric
representation, a transition is symmetry-allowed for the declared
operator.  This is a necessary symmetry condition, not a guarantee
that the matrix element is numerically nonzero: additional quantum
numbers, cancellation, population, or experimental geometry may
suppress the line.

\section{Approximation and identifiability ledger}

A reproducible quantum-chemical claim separates at least the
following errors:
\begin{enumerate}[leftmargin=*,itemsep=2pt]
\item \textbf{physical-model error:} nonrelativistic Hamiltonian,
      nuclear model, environmental and field terms;
\item \textbf{electron-nuclear reduction error:} Born--Oppenheimer,
      retained surfaces, derivative couplings, nuclear quantum
      effects;
\item \textbf{discretization error:} one-particle basis, active space,
      grid, pseudopotential, boundary conditions;
\item \textbf{many-body model error:} determinant ansatz, excitation
      rank, density functional, embedding approximation;
\item \textbf{numerical error:} SCF stationarity, integral and
      diagonalization tolerances, geometry and Hessian convergence;
\item \textbf{property-model error:} response order, transition
      operator, gauge, relativistic picture-change corrections;
\item \textbf{observation error:} instrument response, resolution,
      calibration, noise, estimator, and assignment ambiguity.
\end{enumerate}
Agreement of one final scalar with experiment does not identify these
components separately.

\begin{warning}[Relativistic boundary]
For heavy elements, a relativistic model must state its nuclear charge
distribution, chosen positive-energy or no-pair space, two-electron
terms, and finite-basis kinetic balance.  An unqualified
nonrelativistic Rayleigh--Ritz argument cannot be transferred to the
Dirac operator across its negative-energy continuum.
\end{warning}

\section{Reproducible controls and claim register}

\begin{table}[ht]
\centering
\small
\begin{tabularx}{\linewidth}{L{0.31\linewidth} L{0.23\linewidth} Y}
\toprule
Control & Exact result & Purpose \\
\midrule
\(\ket{\Psi_\pm}=(\ket{12}\pm\ket{34})/\sqrt2\) &
\(\gamma^{(1)}=\tfrac12I_4\), \(\langle B\rangle=\pm1\) &
Non-injectivity of one-body reduction \\
\(C_{\rm o}\mapsto C_{\rm o}U\) &
\(P'=P\), Slater ray unchanged &
Occupied-orbital gauge \\
Diatomic spring \(k\), masses \(m_1,m_2\) &
\(\omega^2=k(m_1^{-1}+m_2^{-1})\) &
Frequency-square and rigid-mode test \\
Two-determinant matrix
\(\bigl(\begin{smallmatrix}0&V\\V&\Delta\end{smallmatrix}\bigr)\) &
\(E_0=(\Delta-\sqrt{\Delta^2+4V^2})/2\le0\) &
Variational correlation sign \\
Gaussian \(L_\sigma\) &
\(\int L_\sigma(\omega-\omega_0)\dd\omega=1\) &
Line-profile normalization \\
Nested Ritz spaces &
\(E_0^{(K+1)}\le E_0^{(K)}\) &
Finite-basis variational control \\
\bottomrule
\end{tabularx}
\caption{Reference values implemented by the P9 regression suite.}
\end{table}

\begin{table}[ht]
\centering
\small
\begin{tabularx}{\linewidth}{L{0.36\linewidth} L{0.18\linewidth} Y}
\toprule
Claim & Status & Scope \\
\midrule
\(\Tr\gamma^{(1)}=N\) and \(\int n=N\) &
Proposition & Normalized \(N\)-electron state \\
One-body reduction is non-injective &
Proposition & Explicit two-electron four-orbital control \\
Occupied rotations preserve \(P\) &
Proposition & Unitary rotations within occupied space \\
\(\lambda_k=\omega_k^2\) &
Theorem & Harmonic mass-weighted Hessian \\
Five or six rigid zero modes &
Qualified theorem & Isolated linear or nonlinear molecule \\
Density uniquely reconstructs any state &
Rejected & Hohenberg--Kohn is a ground-state theorem \\
Orbitals or bonds are unique observables &
Rejected & Representation and scheme must be declared \\
Energy differences alone determine a spectrum &
Rejected & Probe, populations, line shape, and channel are required \\
New quantum-chemical method &
Rejected & Not claimed \\
\bottomrule
\end{tabularx}
\end{table}

\begin{center}
\fcolorbox{obsgreen}{lightgreen}{%
\parbox{0.89\linewidth}{\small
The strict program-level result is a typed inference graph:
Hamiltonian selection, state preparation, many-body reduction,
representation choice, nuclear reduction, ideal property evaluation,
and measurement are distinct operations.  Structural analogy is
retained only where their domains and invariants are explicit.}}
\end{center}

\section*{Protocol minimum}
\addcontentsline{toc}{section}{Protocol minimum}

A complete record states at least: particle numbers and isotope
masses; Hamiltonian and unit context; frame, geometry convention, and
symmetry; pure or mixed state; electronic and nuclear reduction;
basis, overlap threshold, active space, and frozen sectors; orbital
gauge or localization rule; correlation or density-functional model;
SCF and geometry stationarity criteria; Hessian projection and
zero-mode threshold; property operator and gauge; populations, line
profile, instrument response, noise law, estimator, uncertainty, and
identifiability scope.

\clearpage
\section*{References}
\addcontentsline{toc}{section}{References}
\begingroup
\raggedright
\begin{enumerate}[label={[\arabic*]},leftmargin=*,itemsep=3pt]
\item M. Born and R. Oppenheimer, ``Zur Quantentheorie der
Molekeln,'' \emph{Annalen der Physik} 389, 457--484 (1927),
\href{https://doi.org/10.1002/andp.19273892002}
{doi:10.1002/andp.19273892002}.
\item J. C. Slater, ``The Theory of Complex Spectra,''
\emph{Physical Review} 34, 1293--1322 (1929),
\href{https://doi.org/10.1103/PhysRev.34.1293}
{doi:10.1103/PhysRev.34.1293}.
\item C. C. J. Roothaan, ``New Developments in Molecular Orbital
Theory,'' \emph{Reviews of Modern Physics} 23, 69--89 (1951),
\href{https://doi.org/10.1103/RevModPhys.23.69}
{doi:10.1103/RevModPhys.23.69}.
\item P.-O. Lowdin, ``Quantum Theory of Many-Particle Systems. I.
Physical Interpretations by Means of Density Matrices, Natural
Spin-Orbitals, and Convergence Problems in the Method of
Configurational Interaction,'' \emph{Physical Review} 97, 1474--1489
(1955),
\href{https://doi.org/10.1103/PhysRev.97.1474}
{doi:10.1103/PhysRev.97.1474}.
\item A. J. Coleman, ``Structure of Fermion Density Matrices,''
\emph{Reviews of Modern Physics} 35, 668--686 (1963),
\href{https://doi.org/10.1103/RevModPhys.35.668}
{doi:10.1103/RevModPhys.35.668}.
\item P. Hohenberg and W. Kohn, ``Inhomogeneous Electron Gas,''
\emph{Physical Review} 136, B864--B871 (1964),
\href{https://doi.org/10.1103/PhysRev.136.B864}
{doi:10.1103/PhysRev.136.B864}.
\item W. Kohn and L. J. Sham, ``Self-Consistent Equations Including
Exchange and Correlation Effects,'' \emph{Physical Review} 140,
A1133--A1138 (1965),
\href{https://doi.org/10.1103/PhysRev.140.A1133}
{doi:10.1103/PhysRev.140.A1133}.
\item M. Levy, ``Universal variational functionals of electron
densities, first-order density matrices, and natural spin-orbitals
and solution of the \(v\)-representability problem,''
\emph{Proceedings of the National Academy of Sciences USA} 76,
6062--6065 (1979),
\href{https://doi.org/10.1073/pnas.76.12.6062}
{doi:10.1073/pnas.76.12.6062}.
\item C. A. Mead and D. G. Truhlar, ``On the determination of
Born--Oppenheimer nuclear motion wave functions including
complications due to conical intersections and identical nuclei,''
\emph{Journal of Chemical Physics} 70, 2284--2296 (1979),
\href{https://doi.org/10.1063/1.437734}
{doi:10.1063/1.437734}.
\end{enumerate}
\endgroup

\end{document}
